\section{Additional Tables for the Theory Sections}
\label{app:theory_tables}
\label{app:jobdesign_example}
% The house format sets every float [!t]. On the first page of an appendix, that
% puts the leading table at the top of the page that carries the appendix heading,
% above the heading itself, so the table reads as though it belonged to the
% previous appendix. \suppressfloats[t] keeps the top of this page clear, and the
% leading table drops the ! so that it honours the suppression and starts at the
% top of the next page; the tables after it keep [!t] and follow in order behind it.
\suppressfloats[t]

This appendix collects the three tables that the theory sections cite but do not print.

Table~\ref{tab:notation} collects the notation the theory sections use.
Panel~A covers the step-, task-, and strategy-level objects of Sections~\ref{sec:shortrun} and~\ref{sec:results}, and Panel~B the skill, job, and wage objects that Section~\ref{sec:longrun} adds.

\begin{table}[t]
\centering
\caption{Notation Used in the Theory Sections}
\label{tab:notation}
\footnotesize
\begin{tabularx}{\textwidth}{l X}
\toprule
\multicolumn{2}{l}{\textbf{Panel A: Steps, Tasks, and AI Deployment (Section~\ref{sec:shortrun})}} \\
\midrule
\textbf{Symbol} & \textbf{Description} \\
\midrule
\addlinespace
$\alpha \in (0,1]$ & Quality of the general-purpose AI technology \\
\addlinespace
$s_i$ & Production step $i$ \\
$\mathcal{S} = (s_1, \dotsc, s_m)$ & Sequence of the firm's $m$ production steps \\
$\manualTime{i}$ & Time cost of completing step $i$ manually \\
$\AItime{i}$ & Time cost of verifying the output of one AI attempt at step $i$ \\
$d_i$ & Difficulty of step $i$ for AI \\
$q_i = \alpha^{d_i}$ & Probability that AI completes step $i$ successfully \\
\addlinespace
$T_b$ & Task $b$: a manual step or an AI chain (with a standalone augmented step considered a length-one AI chain) \\
$\mathcal{T} = (T_1, \dotsc, T_n)$ & AI deployment strategy; a partition of the steps into $n$ tasks \\
$\timecost{b}$ & Expected time cost of task $b$: $\manualTime{i}$ if a manual step, $\AItime{r}/\prod_{i=\ell}^{r} q_i$ if an AI chain over steps $\ell, \dotsc, r$ \\
\addlinespace
\midrule
\multicolumn{2}{l}{\textbf{Panel B: Skills, Jobs, Wages, and Job Design (Section~\ref{sec:longrun})}} \\
\midrule
\textbf{Symbol} & \textbf{Description} \\
\midrule
\addlinespace
$\manualSkill{i}$ & Skill required to complete step $i$ manually \\
$\AIskill{i}$ & Skill required to verify the output of one AI attempt at step $i$ \\
$\skillcost{b}$ & Skill required by task $b$: $\manualSkill{i}$ if a manual step, $\AIskill{r}$ if an AI chain augmented at step $r$; an automated step requires none \\
\addlinespace
$J_j$ & Job $j$: a contiguous block of tasks assigned to a single worker \\
$\mathcal{J} = (J_1, \dotsc, J_p)$ & Job design; a partition of the tasks into $p$ jobs \\
\addlinespace
$\handofftime{i}$ & Hand-off time incurred when a worker's final step is $s_i$ \\
$\handofftime{}(J_j)$ & Hand-off time of job $J_j$, equal to $\handofftime{i}$ for its final step $s_i$ \\
\addlinespace
$\text{Wage}_j$ & Wage rate of job $J_j$'s worker, equal to the total skill it requires, $\sum_{T_b \in J_j} \skillcost{b}$ \\
$\text{WageBill}_j$ & Per-unit-output labor cost of job $J_j$: $\bigl(\sum_{T_b \in J_j} \skillcost{b}\bigr)\bigl(\handofftime{}(J_j) + \sum_{T_b \in J_j} \timecost{b}\bigr)$ \\
\bottomrule
\end{tabularx}

\begin{minipage}{1\linewidth}
\begin{spacing}{0.2}
{\fontsize{9.5pt}{10pt}\selectfont Notes: Step-level primitives are indexed by $i$, task-level objects by $b$, and jobs by $j$; superscripts $M$ and $A$ denote the manual and AI modes of executing a step, and $H$ denotes a hand-off between consecutive workers.
Panel~A is all the short-run problem requires: with the whole workflow assigned to one worker at a fixed wage, execution time is the only margin and no notion of skill is needed.
Panel~B enters only once the firm chooses how to divide the workflow among workers, and it extends rather than replaces Panel~A, which the long-run problem continues to use.}
\end{spacing}
\end{minipage}
\end{table}

Section~\ref{sec:longrun.specialization} works through a three-task workflow whose four job designs are drawn in Figure~\ref{fig:job_design}.
Table~\ref{tab:job_design} reports what each of them costs, first with hand-off costs switched off and then with them on.
Reading down either panel, a design's total is the sum over its jobs of the job's total skill times the total time it takes, hand-off included, which is Equation~\eqref{eq:wagebill} applied job by job.

Section~\ref{sec:nonlinear} works through a two-step workflow in which the marginal gain from improving AI quality is non-monotone.
Table~\ref{tab:nonmonotone_costs} lists the five AI strategies available there, the cost of each as a function of AI quality $\alpha$, the marginal benefit of raising $\alpha$ under each, and the range of $\alpha$ over which each is cost-minimizing.
Appendix~\ref{app:nonmonotone} states and proves the general result the example illustrates.

\begin{table}[!t]
    \caption{Role of Hand-off Costs in Organizational Structure}
    \vspace{-0.3cm}
    \centering
    \resizebox{0.9\textwidth}{!}{%
    \begin{tabular}{c}
    \begin{minipage}{\textwidth}
    \centering
    Panel (a): Without Hand-off Costs \\[1ex]
    \begin{tabular}{|
  >{\centering\arraybackslash}m{1.7cm}|
  >{\centering\arraybackslash}m{0.9cm}|
  >{\centering\arraybackslash}m{1.8cm}|
  >{\centering\arraybackslash}m{1.8cm}|
  >{\centering\arraybackslash}m{1.8cm}|
  >{\centering\arraybackslash}m{1.6cm}|
  >{\centering\arraybackslash}m{1.7cm}|
  >{\centering\arraybackslash}m{1.7cm}|
  }
\hline
\textbf{Job Design} & \textbf{Job} & \textbf{Job Tasks} & $\sum\hccost{b}$ & $\sum\timecost{b}$ & \textbf{Job Cost} & \textbf{Total Cost} & \textbf{Optimal Design} \\
\hline
\multirow{3}{*}{[1][2][3]} 
  & 1 & $\{1\}$   & 3 & 1 & 3 & \multirow{3}{*}{9} & \multirow{3}{*}{\checkmark} \\
\cline{2-6}
  & 2 & $\{2\}$   & 1 & 2 & 2 &  &  \\
\cline{2-6}
  & 3 & $\{3\}$   & 2 & 2 & 4 &  &  \\
\hline
\multirow{2}{*}{[1,2][3]}
  & 1 & $\{1,2\}$ & 4 & 3 & 12 & \multirow{2}{*}{16} &  \\
\cline{2-6}
  & 2 & $\{3\}$   & 2 & 2 & 4 &  &  \\
\hline
\multirow{2}{*}{[1][2,3]}
  & 1 & $\{1\}$   & 3 & 1 & 3 & \multirow{2}{*}{15} &  \\
\cline{2-6}
  & 2 & $\{2,3\}$ & 3 & 4 & 12 &  &  \\
\hline
[1,2,3] & 1 & $\{1,2,3\}$ & 6 & 5 & 30 & 30 &  \\
\hline
\end{tabular}
    \end{minipage} \\[15ex]

    \begin{minipage}{\textwidth}
    \centering
    Panel (b): Including Hand-off Costs \\[1ex]
    \begin{tabular}{|
  >{\centering\arraybackslash}m{1.7cm}|
  >{\centering\arraybackslash}m{0.9cm}|
  >{\centering\arraybackslash}m{1.8cm}|
  >{\centering\arraybackslash}m{1.8cm}|
  >{\centering\arraybackslash}m{1.8cm}|
  >{\centering\arraybackslash}m{1.6cm}|
  >{\centering\arraybackslash}m{1.7cm}|
  >{\centering\arraybackslash}m{1.7cm}|
  }
\hline
\textbf{Job Design} & \textbf{Job} & \textbf{Job Tasks} & $\sum\hccost{b}$ & $\handofftime{} + \sum\timecost{b}$ & \textbf{Job Cost} & \textbf{Total Cost} & \textbf{Optimal Design} \\
\hline
\multirow{3}{*}{[1][2][3]} 
  & 1 & $\{1\}$   & 3 & 4 & 12 & \multirow{3}{*}{18.5} &  \\
\cline{2-6}
  & 2 & $\{2\}$   & 1 & 2.5 & 2.5 &  &  \\
\cline{2-6}
  & 3 & $\{3\}$   & 2 & 2 & 4 &  &  \\
\hline
\multirow{2}{*}{[1,2][3]}
  & 1 & $\{1,2\}$ & 4 & 3.5 & 14 & \multirow{2}{*}{18} & \multirow{2}{*}{\checkmark} \\
\cline{2-6}
  & 2 & $\{3\}$   & 2 & 2 & 4 &  &  \\
\hline
\multirow{2}{*}{[1][2,3]}
  & 1 & $\{1\}$   & 3 & 4 & 12 & \multirow{2}{*}{24} &  \\
\cline{2-6}
  & 2 & $\{2,3\}$ & 3 & 4 & 12 &  &  \\
\hline
[1,2,3] & 1 & $\{1,2,3\}$ & 6 & 5 & 30 & 30 &  \\
\hline
\end{tabular}
    \end{minipage}
    \end{tabular}
    }
    \label{tab:job_design}
    \footnotesize
    \begin{minipage}{1\linewidth}
\begin{spacing}{0.2}
{\fontsize{9.5pt}{10pt}\selectfont Notes: This table reports total production costs for each possible job design in an example workflow with a fixed automation strategy and $n = 3$ tasks.
The cost parameters of the tasks are given in Section~\ref{sec:longrun.specialization}.
In Panel (a), in the absence of hand-off frictions, full specialization, corresponding to job design [1][2][3], is cost-minimizing.
Once hand-off costs are taken into account, in Panel (b) the optimal organizational structure changes to job design [1,2][3], which avoids the large coordination cost between tasks 1 and 2 at the expense of employing a higher-cost (more skilled) worker to complete those tasks.}
\end{spacing}
\end{minipage}
\end{table}

\begin{table}[!t]
\centering
\caption{Configurations, Costs, and Marginal Gains in Example~\ref{ex:nonmonotone}}
\label{tab:nonmonotone_costs}
\singlespacing
% \footnotesize like every other table float in the paper. At full size this tabular
% is 519pt wide against a 470pt measure and overhangs the right margin by 49pt, which
% is the document's one long-standing "Overfull \hbox (48.87pt too wide)" warning.
\footnotesize
\setlength{\tabcolsep}{10pt}
\begin{tabular}{lccc}
\toprule
Configuration & Time cost ($C$) & $-\mathrm{d}C/\mathrm{d}\alpha$ & Optimal for \\
\midrule
Both steps manual & $\manualTime{1} + \manualTime{2} = 14$ & $0$ & $\alpha \leq 0.50$ \\
\midrule
Step 1 augmented & $\AItime{1}/q_1 + \manualTime{2} = 3.5\alpha^{-11} + 8$ & $38.5/\alpha^{12}$ & never \\
\midrule
Step 2 augmented & $\manualTime{1} + \AItime{2}/q_2 = 6 + 4/\alpha$ & $4/\alpha^{2}$ & $0.50 \leq \alpha \lesssim 0.92$ \\
\midrule
Both steps augmented & $\AItime{1}/q_1 + \AItime{2}/q_2 = 3.5\alpha^{-11} + 4/\alpha$ & $38.5/\alpha^{12} + 4/\alpha^{2}$ & never \\
\midrule
Steps 1--2 chained & $\AItime{2}/(q_1 q_2) = 4\alpha^{-12}$ & $48/\alpha^{13}$ & $\alpha \gtrsim 0.92$ \\
\bottomrule
\end{tabular}
\par\vspace{\fignotesskip}
\begin{minipage}{1\linewidth}
\begin{spacing}{0.2}
{\fontsize{9.5pt}{10pt}\selectfont Notes: The two-step workflow of Example~\ref{ex:nonmonotone}, the AI-hard and AI-easy steps of Section~\ref{sec:comparative_advantage}, in which $\manualTime{1}=6$, $\manualTime{2}=8$, $\AItime{1}=3.5$, $\AItime{2}=4$, and the AI succeeds at step $i$ with probability $q_i = \alpha^{d_i}$ for $d_1 = 11$ and $d_2 = 1$, so that $q_1 = \alpha^{11}$ and $q_2 = \alpha$.
Costs follow from the task time costs of Section~\ref{sec:model.tasks}, and $-\mathrm{d}C/\mathrm{d}\alpha$ is the marginal benefit of improving AI quality.
The last column gives the range of $\alpha$ over which each configuration is cost-minimizing.
Neither configuration in which step~1 is augmented ever minimizes cost.
Augmenting it on its own costs $3.5\alpha^{-11}$, which is prohibitive while AI is unreliable, and by the time AI is reliable enough for that to be affordable, folding step~1 into a chain is cheaper still: an automated step is never prompted or verified on its own, so its management cost is never paid at all.}
\end{spacing}
\end{minipage}
\end{table}
